\section{Conclusion}
\label{a:sec:conclusion}

The constant $A(p)$ organises the quasi-stationary behaviour of subcritical binary
Galton--Watson branching: the limiting conditional mean is $1/A(p)$, the limiting
conditional variance is \cref{a:eq:condvar}, and both are functions of $p$ and
$A(p)$ alone. This chapter has shown that the constant is well defined
(\cref{a:prop:product}), exactly representable as an infinite product
\cref{a:eq:prodFormula} and as a series \cref{a:eq:Aseries}, confined by
\cref{a:prop:bounds} between $\tfrac12\Ac(p)$ and $2\varepsilon/(1+2\varepsilon)$
and by \cref{a:prop:parity} below $\tfrac12$, asymptotically equal to $2\varepsilon$
near criticality with a relative correction of order $\varepsilon\log(1/\varepsilon)$
(\cref{a:thm:nearcritical}), and computable to arbitrary accuracy with a
certificate by the hybrid scheme \cref{a:eq:Apractical}. For every purpose this
thesis has, $A(p)$ is a usable object.

What it is not is elementary, and the reason is structural rather than a failure of
ingenuity. By \cref{a:prop:AisKoenigs} the constant is twice the basin value at
$\tfrac12$ of the Koenigs linearising function of the logistic map at multiplier
$r=2p$, and by \cref{a:thm:ht} that function is hypertranscendental in its own
variable for every $r$ in the subcritical window --- with no exceptional parameters
to exclude, since differentially algebraic Schröder functions exist only over
repelling fixed points, and the window is entirely attracting. This explains the
failure of every closed-form search recorded in \cref{a:sec:GA}, and it explains the
characteristic form of that failure: an object defined only by the dynamics that
generate it will be approximated well by many unrelated expressions and captured by
none.

What the theorem does not decide has been flagged with equal care in
\cref{a:sec:scope}. Transcendence --- even irrationality --- of the individual
values $A(p_0)$ is open, because the 1993 Becker--Bergweiler values theorem is
confined to equal-degree Böttcher pairs and does not transfer, the degree structure
of the Schröder equation being wrong and Mahler's method running the wrong way in
the attracting case. Whether the map $p\mapsto A(p)$ is elementary, or
differentially algebraic in $p$, is open too, and appears to be unaddressed: the
obvious machinery of parametrised difference Galois theory does not apply because
the parameter derivation does not commute with a substitution operator that itself
moves with the parameter. Eleven rational values nonetheless resist all
inverse-symbolic identification against standard constant bases at over $200$-digit
rigour, which is evidence for the conjecture that no closed form exists under any
reading, and is not a proof of it.

Two problems are therefore left open by this chapter, and they are worth separating
from the modelling questions raised in \ChM.

\begin{enumerate}[leftmargin=*]
\item \emph{The value problem.} Is $A(p_0)$ irrational for some, or every, rational
      $p_0\in(0,\tfrac12)$? A positive answer for a single parameter would be the
      first arithmetic result about this constant. The PSLQ battery of
      \cref{a:sec:pslq} bounds the difficulty from one side without touching the
      question.
\item \emph{The parameter problem.} Is $p\mapsto A(p)$ differentially algebraic in
      $p$? This is the reading under which a ``formula for $A(p)$'' would most
      naturally be sought, and it is precisely the reading \cref{a:thm:ht} does not
      address. A negative answer would close the matter; a positive one would be
      more surprising still.
\end{enumerate}

The exceptional elementary linearisations at $r=2$ and $r=4$, both over repelling
fixed points and both outside the branching window, are derived in
\cref{a:app:integrable} and located inside the Becker--Bergweiler classification
there. They are recorded to show what a solvable case looks like, not because they
bear on $A(p)$.
